\documentclass[10pt,a4paper]{amsart}
\usepackage[margin=19mm]{geometry}
\usepackage{amsmath,amssymb,mathtools}
\usepackage[scr=rsfs,cal=euler]{mathalfa}
\usepackage{xfrac}
\usepackage[unicode,hidelinks,bookmarks=false]{hyperref}
\newcommand{\ie}{i.e.\ }
\newcommand{\cf}{cf.\ }
\newcommand{\resp}{respectively\ }
\newcommand{\Subsection}{\subsection}
\providecommand{\nobreakdash}{\nobreak\hbox{-}}
\setlength{\emergencystretch}{2em}
\multlinegap0pt
\usepackage{amsthm}
\newtheorem{theorem}{Theorem}
\numberwithin{theorem}{section}
\newtheorem{lemma}[theorem]{Lemma}
\newtheorem{proposition}[theorem]{Proposition}
\newtheorem{corollary}[theorem]{Corollary}
\theoremstyle{definition}
\newtheorem{definition}[theorem]{Definition}
\newtheorem{example}[theorem]{Example}
\newtheorem{notation}[theorem]{Notation}
\newtheorem{remark}[theorem]{Remark}
\title[On the rational solutions of generalized Abel equations]{On the rational solutions of generalized Abel equations: the composition condition for all-nonlinear Abel equations (Proposition 3.13)}
\author[L. A. Calderon and I. Ojeda]{L.A. Calderón* and I. Ojeda}
\date{Source: arXiv:2605.10591v2 (CC BY 4.0)}
\begin{document}
\maketitle
\vspace{1em}
\noindent\textit{Source and licence.} On the rational solutions of generalized Abel equations. Version arXiv:2605.10591v2 (CC BY 4.0), distributed under the Creative Commons Attribution 4.0 International licence, http://creativecommons.org/licenses/by/4.0/, as recorded in the arXiv OAI licence metadata for this exact version and shown on the abstract page https://arxiv.org/abs/2605.10591v2. This item is an editorial re-presentation of the paper's composition proposition together with the paper's own proof of it and the definitions that proof uses. A full rights notice is delivered at licenses/arxiv-2605.10591-CC-BY-4.0.txt.
\noindent\textit{Editorial scope.} Counted unit: the paper's composition condition for generalized Abel equations whose terms are all nonlinear, printed as Proposition 3.13 (source0.tex lines 378--384), delivered together with the paper's complete original proof of it (source0.tex lines 386--410) as the single primary proof body. No mathematics is written, repaired or re-derived: statement and proof are the paper's own text line for line, in the paper's own notation ($\Bbbk$, $n_1<\cdots<n_s$, $A_i$, $d=\gcd\{n_i-1\}$, $p$, $c_i$, $W$). Required context retained uncounted, because the counted proof is the paper's own assembly over it: the generalized Abel equation itself (source0.tex lines 70--73) and the two displays the proof invokes, namely the denominator reduction (lines 127--129) and the scaled solution condition (lines 159--161). Every cross-reference inside the retained spans resolves inside the delivered document, and no citation command occurs in any retained span, so the delivered document carries no bibliography and no reference or citation command of the delivered text was rewritten or adapted. Printed numbers are the paper's own: the wrapper restores the paper's global equation counter and its section-relative theorem counter before the retained material, so the retained displays carry the paper's own equation numbers (1), (2) and (5), and the counted statement prints as Proposition 3.13. Omitted from this package and not redistributed: the paper's preamble and front matter as such (of which only the title and the author names are reproduced in the wrapper); the introduction and literature review; the Newton-diagram machinery, the simplicity and nonresonance conditions and the remaining results of the paper; the examples; and the closing sections. The paper's own bibliography listing is likewise not redistributed, and since no retained span cites anything, the delivered document carries no bibliography. Printed slips kept literal, among them the paper's own wording and its own choice of symbols. Layout and class: the paper's own class is the standard AMS amsart class with the 12pt and a4paper options, which is not a publisher class; the delivered wrapper is the lane's pinned amsart editorial wrapper at 10pt on A4, to which this package adds the paper's own theorem-like declarations with their shared section-relative counter and its own shorthand macros, all copied verbatim from the paper's source. No publisher class file, no publisher style file and no upstream script is redistributed. Assets: no retained span of this package contains a figure environment or a graphics-inclusion command, no image or artwork file exists in or is redistributed with this package, and the manifest and the delivery evidence both carry an empty asset-dependency list. A full rights notice is delivered at licenses/arxiv-2605.10591-CC-BY-4.0.txt.
	\begin{equation}\label{ecu:main}
		x'=A_s(t)x^{n_s}+\dots+A_1(t)x^{n_1},
		\qquad 1\leq n_1<\cdots<n_s,
	\end{equation}

		\begin{equation}\label{eq:denominator}
			p^{n_s-2}p'+\sum_{i=1}^s A_i p^{n_s-n_i}=0.
		\end{equation}

\setcounter{equation}{4}
		\begin{equation}\label{eq:scaled-condition}
			\sum_{i=1}^s\bigl(\lambda^{n_i-1}-1\bigr)A_i p^{n_s-n_i}=0.
		\end{equation}

\setcounter{section}{3}
\setcounter{theorem}{12}
	\begin{proposition}\label{prop:composition}
		Assume that $1<n_1<\cdots<n_s$, and set $d:=\gcd\{n_1-1,\ldots,n_s-1\}$. If equation~\eqref{ecu:main} has $n_s-1$ distinct nonconstant rational solutions belonging to the same proportionality class, then there exist a nonconstant polynomial $p\in\Bbbk[t]$ and constants $c_i\in\Bbbk^\times$, $i=1,\ldots,s$, such that
		\[
		A_i(t)=c_i p(t)^{n_i-2}p'(t),\qquad i=1,\ldots,s.
		\]
		Consequently, equation~\eqref{ecu:main} satisfies the composition condition with $W(t)=p(t)^d$. In particular, over $\mathbb{R}$, $x=0$ is a composition center on every interval $[a,b]$ such that $p(a)^d=p(b)^d$.
	\end{proposition}

	\begin{proof}
		Let $x=1/p$ be one of the rational solutions. Since the $n_s-1$ solutions belong to the same proportionality class, they can be written as $x_j(t)=\lambda_j/p(t)$, $j=1,\ldots,n_s-1$, for distinct $\lambda_j\in\Bbbk^\times$.
		
		By~\eqref{eq:scaled-condition}, the $\lambda_j$ are roots of $
		H(\lambda):=\sum_{i=1}^s A_i p^{n_s-n_i}\bigl(\lambda^{n_i-1}-1\bigr)$. Since $H$ has degree $n_s-1$ and leading coefficient $A_s$, we have $H(\lambda)=A_s\prod_{j=1}^{n_s-1}(\lambda-\lambda_j)$. Write
		\[
		\prod_{j=1}^{n_s-1}(\lambda-\lambda_j)=\sum_{k=0}^{n_s-1}q_k\lambda^k,\qquad q_{n_s-1}=1.
		\]
		Since every $\lambda_j$ is nonzero, $q_0\neq0$.
		
		On the other hand, since $x=1/p$ is a solution, equation~\eqref{eq:denominator} gives $p^{n_s-2}p'+\sum_{i=1}^s A_i p^{n_s-n_i}=0$.
		Hence $H(0)=-\sum_{i=1}^s A_i p^{n_s-n_i}=p^{n_s-2}p'$. Comparing this with $H(0)=q_0A_s$, we obtain $A_s=q_0^{-1}p^{n_s-2}p'$. Comparison of the coefficient of $\lambda^{n_i-1}$ then yields $A_i p^{n_s-n_i}=q_{n_i-1}A_s$, and therefore
		\[
		A_i=\frac{q_{n_i-1}}{q_0}p^{n_i-2}p',\qquad i=1,\ldots,s.
		\]
		Thus the first assertion follows with $c_i=q_{n_i-1}/q_0$.
		
		Now set $W=p^d$. Since $d\mid n_i-1$ for every $i$, we have $W'=dp^{d-1}p'$ and
		\[
		A_i(t)=\frac{c_i}{d}W(t)^{(n_i-1)/d-1}W'(t).
		\]
		Thus $A_i(t)=\widetilde A_i'(W(t))W'(t)$, where $\widetilde A_i(z)=c_i z^{(n_i-1)/d}/(n_i-1)$, and hence equation~\eqref{ecu:main} satisfies the composition condition with $W=p^d$.
		
		If $\Bbbk=\mathbb{R}$ and $p(a)^d=p(b)^d$, then $W(a)=W(b)$, and hence the return map on $[a,b]$ is the identity in a neighborhood of $x=0$. Therefore $x=0$ is a composition center.
	\end{proof}

\end{document}
